\documentclass[11pt,letterpaper]{article}
\usepackage[margin=1in]{geometry}
\usepackage{graphicx}
\usepackage{booktabs}
\usepackage{hyperref}
\usepackage{fancyhdr}
\usepackage{longtable}

% Header and footer setup
\pagestyle{fancy}
\fancyhf{}
\lhead{\textbf{Westchester County Data Platform}}
\rhead{Reporting Strategy}
\renewcommand{\headrulewidth}{0.4pt}

\title{\textbf{Westchester County Data Platform}\\Reporting Strategy \& Output Catalog}
\author{Arcanum Research Initiative}
\date{October 2025}

\begin{document}

\maketitle

\tableofcontents
\newpage

\section{Overview}

This document catalogs all outputs from the Westchester County Data Platform project and provides guidance on accessing, using, and maintaining the various deliverables. The platform follows Druck organizational standards with a two-folder structure (Output/ and Technical/) ensuring clear separation between user-facing deliverables and implementation details.

\section{Output Structure}

\subsection{Directory Organization}
\begin{verbatim}
Westchester/
├── Output/                    # User-facing deliverables
│   ├── Data/
│   │   ├── Results/          # Processed Excel files
│   │   ├── Source/           # Original data documentation
│   │   └── Robin/            # Robin database integration
│   ├── PDFs/                 # LaTeX-generated reports
│   ├── Charts/               # Static visualizations
│   └── Documentation/        # Analysis guides
└── Technical/                 # Implementation files
    ├── src/                  # Source code
    ├── data/                 # Raw and processed data
    ├── notebooks/            # Jupyter analysis
    ├── docs/                 # LaTeX sources
    └── scripts/              # Automation scripts
\end{verbatim}

\section{Excel Data Files}

\subsection{Druck Compliance Standards}
All Excel files follow mandatory Druck standards:
\begin{itemize}
    \item \textbf{ONE sheet per file} (no multi-tab workbooks)
    \item Machine-readable column names
    \item Professional black \& white formatting
    \item Timestamped filenames: \texttt{[YYYY.MM.DD] description.xlsx}
\end{itemize}

\subsection{Excel File Catalog}

\begin{longtable}{p{7cm}p{8cm}}
\toprule
\textbf{File Name} & \textbf{Description} \\
\midrule
\endhead

\texttt{[2025.10.13] westchester\_county\_demographics.xlsx} & 
County-level demographic statistics from 2022 ACS 5-Year Estimates. Includes population, income, housing, and employment data. \\

\texttt{[2025.10.13] westchester\_census\_tracts.xlsx} & 
Demographic data for all 241 census tracts in Westchester County. Tract-level population, income, and housing metrics. \\

\texttt{[2025.10.13] westchester\_metro\_north\_stations.xlsx} & 
Directory of 56 Metro-North Railroad stations in Westchester County with coordinates and accessibility information. \\

\texttt{[2025.10.13] westchester\_transit\_accessibility.xlsx} & 
Transit accessibility analysis including station counts, ADA compliance rates, and coverage metrics. \\

\bottomrule
\caption{Excel Data Files Inventory}
\end{longtable}

\subsection{How to Use Excel Files}
\begin{enumerate}
    \item Navigate to \texttt{Output/Data/Results/}
    \item Open files in Excel, Google Sheets, or any spreadsheet application
    \item Each file contains exactly ONE sheet with all data
    \item Column headers are self-explanatory and machine-readable
    \item All files are timestamped for version control
\end{enumerate}

\section{PDF Reports}

\subsection{Report Catalog}

\begin{longtable}{p{6cm}p{9cm}}
\toprule
\textbf{Report} & \textbf{Content} \\
\midrule
\endhead

\texttt{westchester\_methodology\_report.pdf} & 
Comprehensive documentation of data sources, processing methods, analytical framework, and technical implementation details. \\

\texttt{westchester\_executive\_summary.pdf} & 
High-level overview of key findings, deliverables, and applications. Non-technical summary for stakeholders and decision-makers. \\

\texttt{westchester\_reporting\_strategy.pdf} & 
This document. Complete catalog of all outputs with usage instructions and maintenance guidelines. \\

\bottomrule
\caption{PDF Reports Inventory}
\end{longtable}

\subsection{Compiling LaTeX Reports}
To regenerate PDF reports from source:
\begin{verbatim}
cd Technical/docs/
pdflatex westchester_methodology_report.tex
pdflatex westchester_executive_summary.tex
pdflatex westchester_reporting_strategy.tex

# Move PDFs to Output/PDFs/
mv *.pdf ../../Output/PDFs/
\end{verbatim}

\section{Web Platform}

\subsection{Starting the Platform}

\textbf{Quick Start:}
\begin{verbatim}
# Backend API
cd Technical/src/api
python main.py
# Runs on http://localhost:8000

# Frontend Application
cd Technical/src/frontend
npm run dev
# Runs on http://localhost:5173
\end{verbatim}

\subsection{Dashboard Catalog}

\begin{longtable}{p{5cm}p{10cm}}
\toprule
\textbf{Dashboard} & \textbf{URL \& Features} \\
\midrule
\endhead

Overview Dashboard & 
\texttt{/dashboards/overview} \\
& High-level county metrics, data availability status, interactive map \\

Demographics Dashboard & 
\texttt{/dashboards/demographics} \\
& Population charts, income distribution, race/ethnicity breakdown, housing statistics \\

Transit Dashboard & 
\texttt{/dashboards/transit} \\
& Station directory, accessibility analysis, coverage maps, 1-mile radius visualization \\

Property Tax Dashboard & 
\texttt{/dashboards/property-tax} \\
& Assessment trends, tax rates, geographic distribution by municipality \\

Budget Dashboard & 
\texttt{/dashboards/budget} \\
& County spending by department, multi-year budget trends, spending priorities \\

Municipal Services Dashboard & 
\texttt{/dashboards/services} \\
& Police, fire, library coverage; emergency response times; facility inventory \\

Municipality Comparison & 
\texttt{/dashboards/comparison} \\
& Side-by-side comparison of demographics, income, and housing across municipalities \\

\bottomrule
\caption{Web Dashboard Catalog}
\end{longtable}

\section{API Endpoints}

\subsection{API Documentation}
Interactive API documentation available at: \texttt{http://localhost:8000/docs}

\subsection{Endpoint Catalog}

\begin{longtable}{p{5cm}p{10cm}}
\toprule
\textbf{Endpoint} & \textbf{Description} \\
\midrule
\endhead

\texttt{GET /} & 
API root with welcome message and status \\

\texttt{GET /api/stats} & 
Platform statistics and data availability \\

\texttt{GET /api/municipalities} & 
List of Westchester County municipalities \\

\texttt{GET /api/transit/stations} & 
Metro-North station data with coordinates \\

\texttt{GET /api/demographics/county} & 
County-level demographics (2022 ACS) \\

\texttt{GET /api/demographics/tracts} & 
Census tract demographics (241 tracts) \\

\texttt{GET /api/demographics/municipalities} & 
Municipality-level demographics \\

\bottomrule
\caption{API Endpoints}
\end{longtable}

\section{Data Update Schedule}

\subsection{Recommended Update Frequency}

\begin{table}[h]
\centering
\begin{tabular}{lll}
\toprule
\textbf{Data Source} & \textbf{Update Frequency} & \textbf{Command} \\
\midrule
Census Demographics & Annually & \texttt{python scripts/download\_all\_data.py} \\
Metro-North GTFS & Monthly & \texttt{python scripts/download\_all\_data.py} \\
NY State Data & Quarterly & \texttt{python scripts/download\_all\_data.py} \\
Excel Files & After data updates & \texttt{python src/processors/generate\_all\_excel.py} \\
\bottomrule
\end{tabular}
\end{table}

\subsection{Update Process}
\begin{enumerate}
    \item Run data download script
    \item Verify data quality
    \item Regenerate Excel files
    \item Validate Druck compliance
    \item Update documentation if needed
    \item Restart web platform
\end{enumerate}

\section{Quality Assurance}

\subsection{Validation Procedures}

\textbf{Excel Validation:}
\begin{verbatim}
python Technical/scripts/validate_excel.py
\end{verbatim}
Checks for ONE sheet per file, validates column names, verifies formatting

\textbf{API Testing:}
\begin{verbatim}
# Access API documentation
http://localhost:8000/docs

# Test each endpoint interactively
\end{verbatim}

\subsection{Compliance Checklist}
\begin{itemize}
    \item[$\square$] All Excel files have exactly ONE sheet
    \item[$\square$] Filenames include timestamp [YYYY.MM.DD]
    \item[$\square$] Column names are machine-readable
    \item[$\square$] All PDFs generated from LaTeX sources
    \item[$\square$] API endpoints return valid JSON
    \item[$\square$] Frontend dashboards load without errors
    \item[$\square$] Documentation is up to date
\end{itemize}

\section{User Support}

\subsection{Documentation Resources}
\begin{itemize}
    \item \textbf{Project README:} \texttt{Projects/Westchester/README.md}
    \item \textbf{Technical Guide:} \texttt{Projects/Westchester/Technical/README.md}
    \item \textbf{Methodology Report:} \texttt{Output/PDFs/westchester\_methodology\_report.pdf}
    \item \textbf{API Docs:} \texttt{http://localhost:8000/docs}
\end{itemize}

\subsection{Common Tasks}

\textbf{Accessing Data:}
\begin{itemize}
    \item Excel files: \texttt{Output/Data/Results/}
    \item PDF reports: \texttt{Output/PDFs/}
    \item Raw data: \texttt{Technical/data/raw/}
\end{itemize}

\textbf{Running Analysis:}
\begin{itemize}
    \item Jupyter notebooks: \texttt{Technical/notebooks/}
    \item Python scripts: \texttt{Technical/scripts/}
    \item API queries: Use interactive docs at \texttt{/docs}
\end{itemize}

\section{Maintenance}

\subsection{Regular Maintenance Tasks}
\begin{enumerate}
    \item \textbf{Monthly:} Update Metro-North GTFS data
    \item \textbf{Quarterly:} Refresh NY State Open Data
    \item \textbf{Annually:} Update Census ACS data when new year is released
    \item \textbf{As Needed:} Add new data sources, expand dashboards
\end{enumerate}

\subsection{Version Control}
\begin{itemize}
    \item All Excel files are timestamped
    \item Git repository tracks all code changes
    \item Documentation updated with each major release
\end{itemize}

\section{Contact Information}

For questions, issues, or requests for additional analysis:
\begin{itemize}
    \item \textbf{Project:} Westchester County Data Platform
    \item \textbf{Organization:} Arcanum Research Initiative
    \item \textbf{Documentation:} See HANDOFF\_DOCUMENTATION.md
\end{itemize}

\end{document}

